\chapter{Local Fields and Class Field Theory}\label{local-fields}
\chapterstatus{draft}

The $p$-adic numbers of Chapter~\ref{p-adic-numbers} are the first example of a \emph{local field}: a field complete with respect to a discrete valuation with finite residue field. Their finite
extensions are local fields too, and they arise as completions of number fields at prime ideals; questions about a number field $K$ near a prime $\mathfrak{p}$ become questions about the local field
$K_\mathfrak{p}$, where Hensel's lemma makes them much easier. This chapter develops the structure of extensions of local fields (ramification index and residue degree, unramified and totally ramified
extensions, Krasner's lemma) and the relation with the decomposition of primes in Chapter~\ref{algebraic-number-theory}. It then states the main theorems of \emph{class field theory}, which describe
all abelian extensions of local and global fields in terms of their arithmetic, with the examples that motivate them: cyclotomic fields, Hilbert class fields and the primes of the form $x^2 + ny^2$.

\noindent\textbf{Prerequisites.} Chapter~\ref{p-adic-numbers} (p-adic Numbers), Chapter~\ref{algebraic-number-theory} (Algebraic Number Theory) and Chapter~\ref{homological-algebra} (Homological Algebra).

\section{Local fields}\label{local-fields:section-local-fields}

\begin{definition}\label{local-fields:definition-local-field}
A \emph{discrete valuation} on a field $K$ is a surjective map $v \colon K \to \ZZ \cup \{\infty\}$ with $v(x) = \infty$ if and only if $x = 0$, $v(xy) = v(x) + v(y)$ and
$v(x + y) \geq \min(v(x), v(y))$. Its \emph{valuation ring} is $\mathcal{O} = \{v \geq 0\}$, with maximal ideal $\mathfrak{p} = \{v > 0\}$ and \emph{residue field} $k = \mathcal{O}/\mathfrak{p}$; a
\emph{uniformiser} is an element $\pi$ with $v(\pi) = 1$. A (non-archimedean) \emph{local field} is a field $K$ with a discrete valuation, complete for the absolute value $|x| = q^{-v(x)}$, whose residue field is
finite with $q$ elements.
\end{definition}

\begin{notation}\label{local-fields:notation-local}
For a local field $K$ we write $\mathcal{O}_K$, $\mathfrak{p}_K$, $k_K$, $\pi_K$, $v_K$ and $q_K$ (or drop the index), and $|\cdot|$ for the normalised absolute value $q^{-v}$. $\mathcal{O}$ is a DVR
(Definition~\ref{commutative-algebra:definition-regular}), every nonzero element of $K$ is $u\pi^m$ with $u \in \mathcal{O}^\times$ and $m = v(x)$, and $\mathfrak{p}^m = \pi^m\mathcal{O}$. A bar denotes reduction
modulo $\mathfrak{p}$.
\end{notation}

\begin{example}\label{local-fields:example-local-fields}
\begin{enumerate}
\item $\QQ_p$ with $v = v_p$, $\mathcal{O} = \ZZ_p$, $\pi = p$, $k = \mathbb{F}_p$ (Chapter~\ref{p-adic-numbers}).
\item The field $\mathbb{F}_q((t))$ of Laurent series $\sum_{i \geq m}a_it^i$, with $v$ the order of vanishing, $\mathcal{O} = \mathbb{F}_q[[t]]$, $\pi = t$, $k = \mathbb{F}_q$. It has characteristic $p$; it is the
completion of the function field $\mathbb{F}_q(t)$ at the place $t = 0$.
\item Finite extensions of $\QQ_p$ (Theorem~\ref{local-fields:theorem-extension-valuation}). It is a theorem that every local field of characteristic $0$ is one of these, and every local field of
characteristic $p$ is isomorphic to some $\mathbb{F}_q((t))$ \cite[Chapter II, Proposition 5.2]{neukirch-ant}.
\end{enumerate}
\end{example}

\begin{proposition}\label{local-fields:proposition-local-structure}
Let $K$ be a local field and $R \subseteq \mathcal{O}$ a set of representatives of $k$ containing $0$.
\begin{enumerate}
\item Every $x \in K^\times$ has a unique expansion $x = \sum_{i \geq m}a_i\pi^i$ with $a_i \in R$, $a_m \neq 0$, $m = v(x)$.
\item $\mathcal{O}$ is compact, $K$ is locally compact and totally disconnected, and $\mathcal{O} \cong \varprojlim\mathcal{O}/\mathfrak{p}^n$.
\item \emph{(Hensel's lemma.)} If $f \in \mathcal{O}[x]$ and $a \in \mathcal{O}$ satisfy $|f(a)| < |f^{\prime}(a)|^2$, then $f$ has a unique root $\alpha \in \mathcal{O}$ with $|\alpha - a| < |f^{\prime}(a)|$.
\item $\mathcal{O}^\times \cong \mu_{q-1} \times (1 + \mathfrak{p})$, where $\mu_{q-1}$ is the group of $(q - 1)$-th roots of unity, which $x \mapsto \bar x$ maps isomorphically onto $k^\times$; and
$K^\times \cong \pi^{\ZZ} \times \mu_{q-1} \times (1 + \mathfrak{p})$.
\end{enumerate}
\end{proposition}

\begin{proof}[Sketch of proof]
(1)--(3) are proved exactly as for $\QQ_p$ (Theorems~\ref{p-adic-numbers:theorem-expansions}, \ref{p-adic-numbers:theorem-topology}, \ref{p-adic-numbers:theorem-inverse-limit}
and~\ref{p-adic-numbers:theorem-hensel}), with $p$ replaced by $\pi$ and the digits $0, \dots, p - 1$ by $R$; the verifications only use the ultrametric inequality, completeness and the finiteness of
$k$. (4) The polynomial $x^{q-1} - 1$ splits into distinct linear factors over $k$ (Theorem~\ref{fields:theorem-finite-fields}), with simple roots, so by (3) each root lifts uniquely to a root in $\mathcal{O}$;
the lifts form $\mu_{q-1}$ and map onto $k^\times$. Every unit is $\omega(u) \cdot u/\omega(u)$ with $\omega(u) \in \mu_{q-1}$ the lift of $\bar u$ and $u/\omega(u) \in 1 + \mathfrak{p}$, uniquely.
\end{proof}

\subsection{Extensions}\label{local-fields:subsection-extensions}

\begin{theorem}\label{local-fields:theorem-extension-valuation}
Let $K$ be a local field and $L/K$ a finite extension of degree $n$. Then $|x|_L = |N_{L/K}(x)|^{1/n}$ is the unique absolute value on $L$ extending $|\cdot|$ on $K$. With it, $L$ is a local field, its
valuation ring $\mathcal{O}_L$ is the integral closure of $\mathcal{O}_K$ in $L$, and $\sigma(x)$ has the same absolute value as $x$ for every $K$-embedding $\sigma$ of $L$ into an extension.
\end{theorem}

\begin{proof}[Sketch of proof]
\emph{Uniqueness.} Two absolute values extending $|\cdot|$ are norms on the finite-dimensional $K$-vector space $L$; since $K$ is complete and locally compact, all norms on $K^n$ are equivalent (the unit
sphere of the max-norm is compact, and a norm is continuous on it and bounded below), so the two absolute values define the same topology and are equal by
Proposition~\ref{p-adic-numbers:proposition-equivalent}, as they agree on $K$.

\emph{Existence.} Multiplicativity of $|N(x)|^{1/n}$ is clear. The ultrametric inequality reduces to: $|N(x)| \leq 1$ implies $|N(1 + x)| \leq 1$. This follows from the statement that $x$ is integral over
$\mathcal{O}_K$ if and only if $N_{L/K}(x) \in \mathcal{O}_K$, which in turn rests on Hensel's lemma in its factorisation form: an irreducible polynomial over $K$ whose constant and leading coefficients are
integral has all coefficients integral \cite[Chapter II, Theorem 4.8]{neukirch-ant}. The same statement identifies $\mathcal{O}_L$ with the integral closure. The value group is a nontrivial subgroup of
$q^{\frac1n\ZZ}$, hence discrete; completeness follows from the equivalence with the max-norm; the residue field is a finite extension of $k$, of degree at most $n$, because lifts to $\mathcal{O}_L$ of elements of $k_L$ linearly independent over $k$ are linearly independent over $K$ (in a relation
$\sum a_i\omega_i = 0$ with $\max|a_i| = 1$, reduction gives a nontrivial relation among the residues), hence finite.
\emph{Conjugates.} $|\sigma(\cdot)|$ is again an absolute value extending $|\cdot|$ on $K(x)$, so it equals $|\cdot|$ by uniqueness.
\end{proof}

\begin{definition}\label{local-fields:definition-e-f}
For a finite extension $L/K$ of local fields, the \emph{ramification index} $e = e(L/K)$ is defined by $\pi_K\mathcal{O}_L = \mathfrak{p}_L^e$, that is, $e = v_L(\pi_K)$, and the \emph{residue degree} is
$f = f(L/K) = [k_L : k_K]$. $L/K$ is \emph{unramified} if $e = 1$ and \emph{totally ramified} if $f = 1$.
\end{definition}

\begin{theorem}\label{local-fields:theorem-ef}
Let $L/K$ be an extension of local fields of degree $n$. Let $\omega_1, \dots, \omega_f \in \mathcal{O}_L$ lift a basis of $k_L$ over $k_K$ (here $f$ is a priori at most $n$). Then the elements
$\omega_i\pi_L^j$, $1 \leq i \leq f$, $0 \leq j < e$, form a basis of $\mathcal{O}_L$ over $\mathcal{O}_K$. In particular $n = ef$.
\end{theorem}

\begin{proof}
\emph{Step 1: a norm estimate.} For $a_1, \dots, a_f \in K$, $|\sum_ia_i\omega_i| = \max_i|a_i|$: dividing by an $a_{i_0}$ of maximal absolute value, the sum becomes an element of $\mathcal{O}_L$ whose reduction is a
nontrivial combination of the $\bar\omega_i$, hence nonzero. So $v_L(\sum_ia_i\omega_i) = e\min_iv_K(a_i)$, a multiple of $e$.

\emph{Step 2: independence.} If $\sum_{i,j}a_{ij}\omega_i\pi_L^j = 0$ with not all $a_{ij} = 0$, then the terms $s_j = \sum_ia_{ij}\omega_i\pi_L^j$ that are nonzero have valuations $\equiv j \pmod e$, pairwise
distinct, so the valuation of the sum is their minimum, which is finite, a contradiction. Moreover $|\sum a_{ij}\omega_i\pi_L^j| = \max_{i,j}|a_{ij}||\pi_L|^j \geq |\pi_K|\max|a_{ij}|$.

\emph{Step 3: spanning.} Let $M$ be the $\mathcal{O}_K$-span of the $\omega_i\pi_L^j$. The quotients $\mathfrak{p}_L^j/\mathfrak{p}_L^{j+1}$, $0 \leq j < e$, are spanned over $k_K$ by the images of the
$\omega_i\pi_L^j$, so $\mathcal{O}_L = M + \mathfrak{p}_L^e = M + \pi_K\mathcal{O}_L$, and by induction $\mathcal{O}_L = M + \pi_K^N\mathcal{O}_L$ for all $N$. Given $x \in \mathcal{O}_L$, choose $m_N \in M$ with
$x - m_N \in \pi_K^N\mathcal{O}_L$. By the estimate of Step 2 the coefficients of $m_N$ form Cauchy sequences in $\mathcal{O}_K$; their limits give $m \in M$ with $x = m$.

\emph{Step 4.} The basis has $ef$ elements and is also a $K$-basis of $L$ (every element of $L$ is an element of $\mathcal{O}_L$ divided by a power of $\pi_K$), so $n = ef$.
\end{proof}

\begin{theorem}[Unramified extensions]\label{local-fields:theorem-unramified}
Let $K$ be a local field with $q$ elements in its residue field. For every $f \geq 1$ there is, inside a fixed algebraic closure, exactly one unramified extension $K_f$ of degree $f$, namely
$K_f = K(\zeta)$ with $\zeta$ a primitive $(q^f - 1)$-th root of unity. It is Galois, and reduction gives an isomorphism $\mathrm{Gal}(K_f/K) \cong \mathrm{Gal}(\mathbb{F}_{q^f}/\mathbb{F}_q)$; so
$\mathrm{Gal}(K_f/K)$ is cyclic, generated by the \emph{Frobenius} $\zeta \mapsto \zeta^q$.
\end{theorem}

\begin{proof}
\emph{Step 1: an unramified extension of degree $f$.} Let $\bar g \in k[x]$ be monic irreducible of degree $f$ (it exists, since $\mathbb{F}_{q^f} = \mathbb{F}_q(\bar\eta)$ for a generator $\bar\eta$), lift it to a
monic $g \in \mathcal{O}_K[x]$, irreducible since its reduction is, and let $L = K(\eta)$ with $g(\eta) = 0$. Then $\eta \in \mathcal{O}_L$ (it is integral) and $\bar\eta$ is a root of $\bar g$, so
$k_L \supseteq \mathbb{F}_{q^f}$ and $f(L/K) \geq f = [L : K]$. By Theorem~\ref{local-fields:theorem-ef}, $e = 1$ and $k_L = \mathbb{F}_{q^f}$.

\emph{Step 2: every unramified $L$ of degree $f$ is $K(\zeta)$.} Its residue field is $\mathbb{F}_{q^f}$, in which $x^{q^f - 1} - 1$ has $q^f - 1$ simple roots; by Hensel's lemma
(Proposition~\ref{local-fields:proposition-local-structure}(3)) they lift to roots in $\mathcal{O}_L$, so $L \supseteq K(\zeta)$. Conversely $K(\zeta)$ has residue field containing $\bar\zeta$, of order $q^f - 1$
(distinct roots of $x^{q^f-1} - 1$ have distinct reductions, the reduced polynomial being separable), hence containing $\mathbb{F}_{q^f}$; so $[K(\zeta) : K] \geq f(K(\zeta)/K) \geq f = [L : K]$, and $L = K(\zeta)$.
This shows uniqueness and, with Step 1, existence.

\emph{Step 3: Galois group.} $K_f$ is the splitting field of $x^{q^f - 1} - 1$, hence Galois of degree $f$. An automorphism is determined by its effect on $\zeta$, and reduction is injective on the roots of
unity of order $q^f - 1$, so $\mathrm{Gal}(K_f/K) \to \mathrm{Gal}(\mathbb{F}_{q^f}/\mathbb{F}_q)$ is injective; both groups have order $f$.
\end{proof}

\begin{proposition}[Eisenstein extensions]\label{local-fields:proposition-eisenstein}
Let $f = x^n + a_{n-1}x^{n-1} + \dots + a_0 \in \mathcal{O}_K[x]$ be Eisenstein ($v_K(a_i) \geq 1$, $v_K(a_0) = 1$). Then $f$ is irreducible, and for a root $\alpha$ the extension $K(\alpha)/K$ is totally ramified of
degree $n$, with uniformiser $\alpha$. Conversely, if $L/K$ is totally ramified, then $L = K(\pi_L)$ and the minimal polynomial of $\pi_L$ is Eisenstein.
\end{proposition}

\begin{proof}
\emph{Eisenstein implies totally ramified.} Let $L = K(\alpha)$, $e = e(L/K) \leq [L : K] \leq n$, and $w = v_L$, so $w(a) = e\,v_K(a)$ on $K$. As $\alpha$ is integral, $w(\alpha) \geq 0$; and
$\alpha^n = -\sum a_i\alpha^i \in \mathfrak{p}_L$, so $w(\alpha) > 0$. For $i \geq 1$, $w(a_i\alpha^i) > e = w(a_0)$, so $nw(\alpha) = w(\alpha^n) = w(a_0) = e$ by the ultrametric inequality (the term $a_0$ is strictly
the smallest). With $w(\alpha) \geq 1$ and $e \leq n$ this forces $w(\alpha) = 1$ and $e = n = [L : K]$; so $f$ is irreducible.

\emph{Conversely.} If $e = n$ and $f = 1$, Theorem~\ref{local-fields:theorem-ef} (with $\omega_1 = 1$) shows that $1, \pi_L, \dots, \pi_L^{n-1}$ is a basis, so $L = K(\pi_L)$. Its conjugates have the same absolute
value $|\pi_L| < 1$ (Theorem~\ref{local-fields:theorem-extension-valuation}), so the coefficients of its minimal polynomial, symmetric functions of them, have absolute value $< 1$, and
$|a_0| = |N(\pi_L)| = |\pi_L|^n = |\pi_K|$.
\end{proof}

\begin{example}\label{local-fields:example-cyclotomic-local}
For $n \geq 1$, $\Phi_{p^n}(x) = \Phi_p(x^{p^{n-1}})$ satisfies $\Phi_{p^n}(x + 1) \equiv x^{\varphi(p^n)} \pmod p$ (as $(x + 1)^{p^k} \equiv x^{p^k} + 1$) and $\Phi_{p^n}(1) = p$; so it is Eisenstein, and
$\QQ_p(\zeta_{p^n})$ is totally ramified of degree $\varphi(p^n) = p^{n-1}(p - 1)$ with uniformiser $\zeta_{p^n} - 1$ and Galois group $(\ZZ/p^n)^\times$. On the other hand $\QQ_p(\zeta_m)$ for $p \nmid m$ is
unramified, of degree the order of $p$ modulo $m$, by Theorem~\ref{local-fields:theorem-unramified}. Every extension $L/K$ contains a largest unramified subextension $L_0$, with $[L_0 : K] = f$ and $L/L_0$
totally ramified \cite[Chapter II, \S 7]{neukirch-ant}.
\end{example}

\begin{lemma}[Krasner]\label{local-fields:lemma-krasner}
Let $\alpha$ be separable over $K$ with conjugates $\alpha = \alpha_1, \alpha_2, \dots, \alpha_d$ in an algebraic closure, and let $\beta$ satisfy $|\beta - \alpha| < |\alpha - \alpha_i|$ for all $i \geq 2$. Then
$K(\alpha) \subseteq K(\beta)$.
\end{lemma}

\begin{proof}
Otherwise, in a Galois extension containing $\alpha$ and $\beta$ there is an automorphism $\sigma$ fixing $K(\beta)$ with $\sigma\alpha = \alpha_i \neq \alpha$. By uniqueness of the extended absolute value
(Theorem~\ref{local-fields:theorem-extension-valuation}), $|\beta - \sigma\alpha| = |\sigma(\beta - \alpha)| = |\beta - \alpha|$. Then
$|\alpha - \alpha_i| \leq \max(|\alpha - \beta|, |\beta - \sigma\alpha|) = |\beta - \alpha| < |\alpha - \alpha_i|$, a contradiction.
\end{proof}

\begin{corollary}\label{local-fields:corollary-krasner-polynomials}
Let $f \in K[x]$ be monic, separable and irreducible of degree $d$. Every monic $g \in K[x]$ of degree $d$ whose coefficients are close enough to those of $f$ is irreducible, and a root of $g$ generates
the same field as a suitable root of $f$. Consequently every finite extension of $\QQ_p$ is generated by a root of a polynomial with rational coefficients, so it is $\QQ_p \cdot F$ for a number field $F$.
\end{corollary}

\begin{proof}[Sketch of proof]
The roots of a monic polynomial depend continuously on its coefficients: if $g$ is close to $f$, then $|f(\beta)| = |f(\beta) - g(\beta)|$ is small for every root $\beta$ of $g$, so $\prod_i|\beta - \alpha_i|$ is small and
$\beta$ is close to some root $\alpha$ of $f$, closer than $\min_{i \neq j}|\alpha_i - \alpha_j|$. Krasner's lemma gives $K(\alpha) \subseteq K(\beta)$, and $[K(\beta) : K] \leq d = [K(\alpha) : K]$ gives equality and
the irreducibility of $g$. For the last statement, approximate the minimal polynomial of a generator by one with coefficients in $\QQ$, which is dense in $\QQ_p$.
\end{proof}

\subsection{Completions of number fields}\label{local-fields:subsection-completions}

Let $K$ be a number field and $\mathfrak{p}$ a nonzero prime of $\mathcal{O}_K$ above $p$. Since $(\mathcal{O}_K)_\mathfrak{p}$ is a DVR, $v_\mathfrak{p}(x)$, the exponent of $\mathfrak{p}$ in $x\mathcal{O}_K$, is a discrete
valuation on $K$; the completion $K_\mathfrak{p}$ is a local field with residue field $\mathcal{O}_K/\mathfrak{p}$.

\begin{theorem}\label{local-fields:theorem-completions}
\begin{enumerate}
\item $K_\mathfrak{p}$ is a finite extension of $\QQ_p$ with $e(K_\mathfrak{p}/\QQ_p) = e_\mathfrak{p}$ and $f(K_\mathfrak{p}/\QQ_p) = f_\mathfrak{p}$, the invariants of
Definition~\ref{algebraic-number-theory:definition-ramification}; so $[K_\mathfrak{p} : \QQ_p] = e_\mathfrak{p}f_\mathfrak{p}$.
\item The natural map $K \otimes_\QQ\QQ_p \to \prod_{\mathfrak{p} \mid p}K_\mathfrak{p}$ is an isomorphism; equivalently, if $K = \QQ(\theta)$ with minimal polynomial $F$, the factorisation $F = \prod F_\mathfrak{p}$ into
irreducibles over $\QQ_p$ corresponds to the primes above $p$, with $\deg F_\mathfrak{p} = e_\mathfrak{p}f_\mathfrak{p}$.
\item If $K/\QQ$ is Galois, the decomposition group $D_\mathfrak{p}$ (Definition~\ref{algebraic-number-theory:definition-decomposition-group}) is identified with $\mathrm{Gal}(K_\mathfrak{p}/\QQ_p)$, and the inertia group
with the subgroup fixing the maximal unramified subextension.
\end{enumerate}
\end{theorem}

\begin{proof}[Sketch of proof]
(1) $K_\mathfrak{p}$ is generated over $\QQ_p$ by $K$, hence by $\theta$, so it is finite. Completion changes neither the residue field nor the value group, so $f(K_\mathfrak{p}/\QQ_p) = f_\mathfrak{p}$ and
$v_\mathfrak{p}(p) = e_\mathfrak{p}$ gives $e(K_\mathfrak{p}/\QQ_p) = e_\mathfrak{p}$; then Theorem~\ref{local-fields:theorem-ef}. (2) The map is surjective because $K$ is dense in $\prod_\mathfrak{p}K_\mathfrak{p}$ (weak
approximation, proved from the Chinese remainder theorem as in Lemma~\ref{p-adic-numbers:lemma-approximation}), and both sides have dimension $n = \sum e_\mathfrak{p}f_\mathfrak{p}$ over $\QQ_p$
(Theorem~\ref{algebraic-number-theory:theorem-sum-ef}). $K \otimes \QQ_p = \QQ_p[x]/(F) = \prod\QQ_p[x]/(F_i)$, and each factor is one of the $K_\mathfrak{p}$. (3) Each $\sigma \in D_\mathfrak{p}$ is continuous for
$v_\mathfrak{p}$ and extends to $K_\mathfrak{p}$; this gives an injection $D_\mathfrak{p} \to \mathrm{Gal}(K_\mathfrak{p}/\QQ_p)$ between groups of order $e_\mathfrak{p}f_\mathfrak{p}$. See \cite[Chapter II, \S 8--9]{neukirch-ant}.
\end{proof}

\begin{example}\label{local-fields:example-completions}
In $\QQ(i)$: $5 = (2 + i)(2 - i)$ splits, and correspondingly $x^2 + 1$ has two roots in $\QQ_5$ (Hensel, from $2^2 \equiv -1 \pmod 5$), so $\QQ(i) \otimes \QQ_5 \cong \QQ_5 \times \QQ_5$; $3$ is inert and
$\QQ(i)_{(3)} = \QQ_3(i)$ is the unramified quadratic extension of $\QQ_3$; $2$ ramifies and $\QQ_2(i) = \QQ_2(i - 1)$, with $i - 1$ a root of the Eisenstein polynomial $x^2 + 2x + 2$.
\end{example}

\begin{remark}\label{local-fields:remark-ramification-groups}
The inertia group carries a finer filtration by the \emph{higher ramification groups} $G_i = \{\sigma : v_L(\sigma x - x) \geq i + 1 \text{ for } x \in \mathcal{O}_L\}$; $G_0/G_1$ is cyclic of order prime to
$p$ (tame ramification) and $G_1$ is a $p$-group (wild ramification). The \emph{different} $\mathfrak{D}_{L/K}$, the inverse of the dual of $\mathcal{O}_L$ under the trace form, has
$v_L(\mathfrak{D}) = \sum_{i \geq 0}(|G_i| - 1)$, and its norm is the discriminant; this refines Dedekind's theorem that the ramified primes divide the discriminant \cite[Chapter III, \S 2]{neukirch-ant},
\cite{serre-local-fields}.
\end{remark}

\section{Class field theory}\label{local-fields:section-class-field-theory}

Class field theory describes the abelian extensions of a local or global field $K$, that is, the Galois extensions with abelian group, through objects built from $K$ itself: the multiplicative group
$K^\times$ for local fields, and the idele class group for number fields. Its proofs (by Galois cohomology, or by analytic methods for the global part) are long; here we state the theorems and work out
the examples that make them concrete. Complete proofs are in \cite{neukirch-ant} and \cite{cassels-frohlich}.

\subsection{Local class field theory}\label{local-fields:subsection-local-cft}

\begin{theorem}[Local reciprocity]\label{local-fields:theorem-local-cft}
Let $K$ be a local field. There is a unique homomorphism $\mathrm{rec}_K \colon K^\times \to \mathrm{Gal}(K^{\mathrm{ab}}/K)$ to the Galois group of the maximal abelian extension such that:
\begin{enumerate}
\item for every uniformiser $\pi$ and finite unramified extension $L/K$, $\mathrm{rec}_K(\pi)|_L$ is the Frobenius of Theorem~\ref{local-fields:theorem-unramified};
\item for every finite abelian extension $L/K$, $\mathrm{rec}_K$ induces an isomorphism $K^\times/N_{L/K}(L^\times) \cong \mathrm{Gal}(L/K)$.
\end{enumerate}
Moreover $L \mapsto N_{L/K}(L^\times)$ is an inclusion-reversing bijection between finite abelian extensions and open subgroups of finite index of $K^\times$ (the \emph{existence theorem})
\cite[Chapter V, \S 1]{neukirch-ant}.
\end{theorem}

\begin{example}\label{local-fields:example-local-cft-qp}
For $K = \QQ_p$: $\QQ_p^\times = p^\ZZ \times \ZZ_p^\times$, and $\QQ_p^{\mathrm{ab}}$ is the compositum of the unramified extensions $\QQ_p(\zeta_m)$, $p \nmid m$, and of the totally ramified
$\QQ_p(\zeta_{p^n})$ (\emph{local Kronecker--Weber theorem}). The reciprocity map sends $p$ to the Frobenius on the first and to the identity on the second, and $u \in \ZZ_p^\times$ to the identity on the
first and to $\zeta_{p^n} \mapsto \zeta_{p^n}^{u^{-1}}$ on the second. For the quadratic extensions, $\mathrm{rec}(b)$ acts on $\QQ_p(\sqrt a)$ by the sign $(a, b)_p$: the Hilbert symbol of
Chapter~\ref{p-adic-numbers} is the local reciprocity map for quadratic extensions, which is why $(a, b)_p = 1$ exactly when $b$ is a norm from $\QQ_p(\sqrt a)$
(Proposition~\ref{p-adic-numbers:proposition-hilbert-norm}).
\end{example}

\subsection{Adeles and ideles}\label{local-fields:subsection-adeles}

\begin{definition}\label{local-fields:definition-adeles}
The \emph{adele ring} $\mathbb{A}_K$ of a number field $K$ is the restricted product of the completions $K_v$ over all places $v$ (the primes $\mathfrak{p}$ and the archimedean places, the real and complex
embeddings up to conjugation): the elements $(x_v)$ with $x_v \in \mathcal{O}_{K_v}$ for all but finitely many $v$, topologised so that $\prod_{v \mid \infty}K_v \times \prod_{\mathfrak{p}}\mathcal{O}_{K_\mathfrak{p}}$
is open with the product topology. The \emph{idele group} $\mathbb{I}_K = \mathbb{A}_K^\times$ is the restricted product of the $K_v^\times$ with respect to the $\mathcal{O}_{K_\mathfrak{p}}^\times$. $K$ embeds
diagonally in both, and $C_K = \mathbb{I}_K/K^\times$ is the \emph{idele class group}.
\end{definition}

\begin{proposition}\label{local-fields:proposition-adeles-q}
$\QQ$ is a discrete subgroup of $\mathbb{A}_\QQ$, and $\mathbb{A}_\QQ = \QQ + \bigl([0, 1) \times \prod_p\ZZ_p\bigr)$ with $\QQ \cap \bigl([0, 1) \times \prod_p\ZZ_p\bigr) = \{0\}$; hence $\mathbb{A}_\QQ/\QQ$ is compact. Also
$\mathbb{I}_\QQ = \QQ^\times \times \bigl(\RR_{>0} \times \prod_p\ZZ_p^\times\bigr)$.
\end{proposition}

\begin{proof}
\emph{Decomposition.} Let $x = (x_\infty, x_2, x_3, \dots) \in \mathbb{A}_\QQ$. For the finitely many $p$ with $x_p \notin \ZZ_p$, the $p$-adic expansion gives $r_p \in \ZZ[\frac1p]$ with $x_p - r_p \in \ZZ_p$
(the polar part). Then $r = \sum r_p \in \QQ$ satisfies $x_p - r \in \ZZ_p$ for all $p$, since $r_\ell \in \ZZ_p$ for $\ell \neq p$. Subtracting the integer $\lfloor x_\infty - r\rfloor$ as well, $x - r^{\prime}$ has
real component in $[0, 1)$ and $p$-adic components in $\ZZ_p$. \emph{Uniqueness and discreteness.} A rational number in $\ZZ_p$ for all $p$ is an integer, and an integer in $[0, 1)$ is $0$; the open set
$(-\frac12, \frac12) \times \prod\ZZ_p$ meets $\QQ$ only in $0$. \emph{Compactness.} $\mathbb{A}_\QQ/\QQ$ is the image of the compact set $[0, 1] \times \prod_p\ZZ_p$ (Tychonoff). \emph{Ideles.} An idele $(x_v)$ has
$x_p \in \ZZ_p^\times$ for almost all $p$; with $r = \pm\prod_pp^{v_p(x_p)}$, the sign that of $x_\infty$, $x/r$ has unit components at all $p$ and positive real component, and $r$ is unique.
\end{proof}

So $C_\QQ \cong \RR_{>0} \times \hat{\ZZ}^\times$, where $\hat{\ZZ}^\times = \prod_p\ZZ_p^\times = \varprojlim(\ZZ/m)^\times$; the second factor is the Galois group of $\QQ(\zeta_m, m \geq 1)$ over $\QQ$, which is the content
of the Kronecker--Weber theorem (Theorem~\ref{algebraic-number-theory:theorem-kronecker-weber}) in the language of global class field theory.

\subsection{Global class field theory}\label{local-fields:subsection-global-cft}

\begin{theorem}[Artin reciprocity and the existence theorem]\label{local-fields:theorem-global-cft}
Let $K$ be a number field. There is a continuous surjective homomorphism $\mathrm{rec}_K \colon C_K \to \mathrm{Gal}(K^{\mathrm{ab}}/K)$ whose restriction to each $K_v^\times$ is the local reciprocity map of
$K_v$ (composed with the inclusion of decomposition groups). For every finite abelian $L/K$ it induces $C_K/N_{L/K}C_L \cong \mathrm{Gal}(L/K)$, and $L \mapsto N_{L/K}C_L$ is a bijection between finite abelian
extensions and open subgroups of finite index of $C_K$. For an unramified prime $\mathfrak{p}$, the idele $(1, \dots, \pi_\mathfrak{p}, \dots, 1)$ maps to the Frobenius $\mathrm{Frob}_\mathfrak{p}$
\cite[Chapter VI, \S\S 5--6]{neukirch-ant}.
\end{theorem}

In terms of ideals: for a modulus $\mathfrak{m}$, the \emph{ray class group} $\mathrm{Cl}_\mathfrak{m}$ is the group of fractional ideals prime to $\mathfrak{m}$ modulo principal ideals $(\alpha)$ with
$\alpha \equiv 1 \pmod{\mathfrak{m}}$ (and positive at the real places dividing $\mathfrak{m}$); for every finite abelian $L/K$ there is $\mathfrak{m}$ such that the \emph{Artin map}
$\mathfrak{p} \mapsto \mathrm{Frob}_\mathfrak{p}$ factors through a surjection $\mathrm{Cl}_\mathfrak{m} \to \mathrm{Gal}(L/K)$. For $K = \QQ$, $L = \QQ(\zeta_m)$ and $\mathfrak{m} = m\infty$, this is the isomorphism
$(\ZZ/m)^\times \cong \mathrm{Gal}(\QQ(\zeta_m)/\QQ)$, $q \mapsto \sigma_q$, proved for prime $m$ in Theorem~\ref{algebraic-number-theory:theorem-cyclotomic-integers}(3).

\begin{theorem}[Hilbert class field]\label{local-fields:theorem-hilbert-class-field}
Every number field $K$ has a maximal unramified (at all places, including the real ones) abelian extension $H$, the \emph{Hilbert class field}. The Artin map gives $\mathrm{Cl}_K \cong \mathrm{Gal}(H/K)$, so
$[H : K] = h_K$, and a prime $\mathfrak{p}$ of $K$ splits completely in $H$ if and only if it is principal \cite[Chapter VI, \S 6]{neukirch-ant}.
\end{theorem}

\begin{example}\label{local-fields:example-x2-5y2}
Let $K = \QQ(\sqrt{-5})$, with $h_K = 2$ (Example~\ref{algebraic-number-theory:example-class-numbers}). The field $L = K(i) = \QQ(i, \sqrt5)$ is an unramified quadratic extension of $K$: only $2$ and $5$
could ramify; $5$ is unramified in $\QQ(i)$, so its ramification index in $L$ is at most $[L : \QQ(i)] = 2$, which it already has in $K$, so nothing new ramifies above $5$; likewise above $2$, using $\QQ(\sqrt5)$;
and $K$ has no real places. By Theorem~\ref{local-fields:theorem-hilbert-class-field}, $L = H$. Now for a prime $p \neq 2, 5$:
\[
\begin{aligned}
p = x^2 + 5y^2 &\iff p\mathcal{O}_K = \mathfrak{p}\bar{\mathfrak{p}} \text{ with } \mathfrak{p} = (x + y\sqrt{-5}) \text{ principal}\\
&\iff p \text{ splits completely in } H\\
&\iff \Bigl(\frac{-1}{p}\Bigr) = \Bigl(\frac{5}{p}\Bigr) = 1 \iff p \equiv 1, 9 \pmod{20}.
\end{aligned}
\]
For instance $29 = 3^2 + 5 \cdot 2^2$ and $41 = 6^2 + 5 \cdot 1^2$, while $3 \cdot 7 = 21 = 1 + 5 \cdot 2^2$ although neither $3$ nor $7$ is of the form: they lie in the non-principal class. This is the
answer to the question raised in Chapter~\ref{number-theory} \cite{cox-primes}.
\end{example}

\begin{remark}\label{local-fields:remark-cft-proofs}
The local theory is proved with Galois cohomology: for a finite Galois extension $L/K$ of local fields, $H^2(\mathrm{Gal}(L/K), L^\times)$ is cyclic of order $[L : K]$ with a canonical generator (the
\emph{fundamental class}), and cup product with it gives Tate's isomorphism $\hat H^{-2}(G, \ZZ) = G^{\mathrm{ab}} \cong \hat H^0(G, L^\times) = K^\times/NL^\times$, whose inverse is $\mathrm{rec}_K$; the Brauer
group $\mathrm{Br}(K) \cong \QQ/\ZZ$ appears here. The global theory adds the cohomology of idele class groups and the analytic or algebraic proof of the second inequality
$[C_K : N_{L/K}C_L] \leq [L : K]$ \cite{cassels-frohlich}. Tate's thesis recasts Hecke's analytic continuation of $L$-functions as harmonic analysis on $\mathbb{A}_K$ and $\mathbb{I}_K$, with the Poisson
summation formula for $\mathbb{A}_K/K$ (Proposition~\ref{local-fields:proposition-adeles-q} is the case $K = \QQ$ of its compactness); the non-abelian generalisation of class field theory is the Langlands
programme.
\end{remark}

\section{Exercises}\label{local-fields:section-exercises}

\begin{exercise}\label{local-fields:exercise-quadratic-extensions}
Let $p$ be odd. Show that $\QQ_p$ has exactly three quadratic extensions, that one of them is unramified, and that the other two are totally ramified.
\end{exercise}

\begin{solution}
Quadratic extensions are $\QQ_p(\sqrt a)$ with $a$ running through the nontrivial classes of $\QQ_p^\times/\QQ_p^{\times2}$, whose representatives are $\varepsilon$, $p$, $\varepsilon p$ with $\varepsilon$ a
non-residue (Corollary~\ref{p-adic-numbers:corollary-square-classes}). $\QQ_p(\sqrt\varepsilon)$ has residue field $\mathbb{F}_p(\sqrt{\bar\varepsilon}) = \mathbb{F}_{p^2}$, so it is unramified (it is
$\QQ_p(\zeta_{p^2 - 1})$ by Theorem~\ref{local-fields:theorem-unramified}). $x^2 - p$ and $x^2 - \varepsilon p$ are Eisenstein, so the other two are totally ramified
(Proposition~\ref{local-fields:proposition-eisenstein}).
\end{solution}

\begin{exercise}\label{local-fields:exercise-q2-extensions}
Determine $e$ and $f$ for $\QQ_2(\sqrt{-3})$, $\QQ_2(\sqrt{-1})$ and $\QQ_2(\sqrt 2)$.
\end{exercise}

\begin{solution}
$\QQ_2(\sqrt{-3}) = \QQ_2(\zeta_3)$ is unramified of degree $2$, as $2$ has order $2$ modulo $3$ (Example~\ref{local-fields:example-cyclotomic-local}): $e = 1$, $f = 2$. $\QQ_2(i) = \QQ_2(i - 1)$ with $i - 1$ a
root of the Eisenstein polynomial $x^2 + 2x + 2$: $e = 2$, $f = 1$. $x^2 - 2$ is Eisenstein: $e = 2$, $f = 1$.
\end{solution}

\begin{exercise}\label{local-fields:exercise-x2-6y2}
Assuming that the Hilbert class field of $K = \QQ(\sqrt{-6})$ is $K(\sqrt{-3}) = \QQ(\sqrt{-6}, \sqrt 2)$, show that a prime $p > 3$ is of the form $x^2 + 6y^2$ if and only if $p \equiv 1, 7 \pmod{24}$.
Check that $K(\sqrt{-3})/K$ is unramified.
\end{exercise}

\begin{solution}
$h_K = 2$ (Exercise~\ref{algebraic-number-theory:exercise-sqrt-minus-six}). \emph{Unramified:} only $2$ and $3$ ramify in $H = \QQ(\sqrt{-6}, \sqrt2)$ over $\QQ$; $3$ is unramified in $\QQ(\sqrt2)$ and $2$ in
$\QQ(\sqrt{-3})$, so as in Example~\ref{local-fields:example-x2-5y2} their ramification indices in $H$ are $2$, already attained in $K$. \emph{Splitting:} $p = x^2 + 6y^2$ iff $p$ splits completely in $H$ iff
$p$ splits in $\QQ(\sqrt{-6})$, $\QQ(\sqrt 2)$ and $\QQ(\sqrt{-3})$; it suffices that it splits in the last two, which means $(\frac2p) = 1$ and $(\frac{-3}{p}) = 1$, that is, $p \equiv \pm1 \pmod 8$ and
$p \equiv 1 \pmod 3$, that is, $p \equiv 1, 7 \pmod{24}$. Examples: $7 = 1 + 6$, $31 = 25 + 6$, $73 = 49 + 24$.
\end{solution}

\begin{exercise}\label{local-fields:exercise-roots-of-unity-qp}
Show that the roots of unity in $\QQ_p$ are exactly $\mu_{p-1}$ for $p$ odd and $\{\pm1\}$ for $p = 2$, using Example~\ref{local-fields:example-cyclotomic-local}.
\end{exercise}

\begin{solution}
Roots of unity of order prime to $p$ in $\QQ_p$: if $\zeta_m \in \QQ_p$ with $p \nmid m$, then $\QQ_p(\zeta_m) = \QQ_p$ has degree $1$, the order of $p$ modulo $m$, so $m \mid p - 1$. Roots of $p$-power
order: $\QQ_p(\zeta_{p^n})$ has degree $\varphi(p^n)$, which is $1$ only for $p^n = 2$. A root of unity of order $m = p^am^{\prime}$ generates roots of orders $p^a$ and $m^{\prime}$. So the group is $\mu_{p-1}$ for odd
$p$, and $\mu_2$ for $p = 2$ (where $p - 1 = 1$).
\end{solution}
